% !TeX root = ../thuthesis-example.tex

\chapter{需求分析（设计目标）}

\section{总体目标}

系统的总体目标是构建一套可运行、可复现、可对比的电商行为大数据分析实验平台。平台既要能处理真实 Kaggle Oct/Nov 行为数据，又要适合本机 Docker Desktop 环境，不能为了追求形式上的“集群”而频繁 OOM 或占满磁盘。最终交付应包含四类成果：第一，ARM64 Docker Hadoop YARN 实验集群；第二，Spark 分析 pipeline 与内存优化实现；第三，1\% 和 5\% benchmark 对照结果；第四，前端 Ops 页面和 LaTeX 报告中的实验图表。

\section{功能需求}

\subsection{数据处理功能}

数据处理链路需要支持从原始 CSV 到清洗缓存的完整闭环。原始数据来自 Kaggle 电商行为数据集，字段包含 \texttt{event\_time}、\texttt{event\_type}、\texttt{product\_id}、\texttt{category\_id}、\texttt{category\_code}、\texttt{brand}、\texttt{price}、\texttt{user\_id} 和 \texttt{user\_session}。系统需要完成字段类型转换、异常价格过滤、重复事件处理、会话事实表生成、日趋势、转化漏斗、商品推荐、关联图谱、异常信号、cohort、forecasting、experimentation 和 feature mart 等分析模块。

\subsection{集群运行功能}

系统需要通过 Docker Compose 启动 Hadoop YARN 实验集群，并支持 Spark submit 在 YARN client mode 下运行。NameNode、ResourceManager 和 Spark History Server 必须提供可访问 UI，便于答辩时展示 HDFS 文件、YARN application 和 Spark event log。由于已有 \texttt{master} 容器占用默认端口，新集群必须使用偏移端口，避免端口冲突。

\subsection{运维展示功能}

前端 Ops 页面需要展示当前 Spark 作业状态、质量门禁、HDFS 输入快照、Spark application id、Spark History 指标状态、最近运行记录和 benchmark 证据。新增 benchmark 展示应覆盖 1\% 五组、5\% 两组实验结果，并突出 YARN-only、AQE、算法优化和 Parquet 输入之间的差异。

\section{非功能需求}

\subsection{架构兼容性}

本项目必须兼容 ARM Mac。所有 Hadoop/Spark 相关镜像必须使用 ARM64 架构，不能下载 amd64 镜像，也不能依赖 \texttt{linux/amd64} 平台模拟。最终镜像不采用 \texttt{docker commit master} 这种粗暴方式，而是构建只包含 Java、Hadoop、必要配置和启动脚本的精简镜像。

\subsection{内存稳定性}

系统目标不是承诺任意数据规模都不会 OOM，而是在本机资源范围内显著降低 OOM 风险，并通过实验给出证据。driver 端不得收集完整明细；大输出应落地为 Parquet；JSON 只保留 summary、quality、manifest 和 preview；模块缓存必须使用 \texttt{MEMORY\_AND\_DISK}，并在模块结束后释放。对于关联分析、推荐和异常检测，应设置候选上界，避免二次级组合膨胀。

\subsection{证据可追溯性}

每次运行应保留 run manifest、配置 hash、input snapshot、quality report、Spark application id 和 output artifact row counts。Benchmark 结果应保存到 \texttt{data/benchmarks}，报告图表保存到 \texttt{报告/figures}。前端展示的数据不应依赖人工复制终端输出，而应由 API 自动读取正式 benchmark 文件。

\section{实验设计}

实验规模分为 1\% 用户级样本和 5\% 用户级样本。1\% 样本用于五组对照，5\% 样本用于压力验证。对照组设计如表\ref{tab:experiment-groups}所示。

\begin{table}[htbp]
\centering
\caption{实验对照组设计}
\label{tab:experiment-groups}
\begin{tabularx}{0.95\textwidth}{lX}
\toprule
\textbf{对照组} & \textbf{实验目的} \\
\midrule
Baseline local CSV & 观察单机 local Spark 在 1\% 样本上的基线耗时和 driver 内存。 \\
YARN-only CSV & 只切换到 YARN client mode，不做算法优化，验证“只加 YARN”是否有效。 \\
YARN + AQE CSV & 启用 AQE、skew join 和合理 shuffle partitions，验证执行计划优化收益。 \\
YARN + algorithm CSV & 在 AQE 基础上限制 collect、关联组合和 preview 输出，验证算法层优化收益。 \\
YARN + Parquet & 使用清洗后的 Parquet 输入，验证列式数据布局收益。 \\
\bottomrule
\end{tabularx}
\end{table}

\section{验收指标}

系统验收指标包括工程指标和实验指标两类。工程指标要求 \texttt{docker compose --profile yarn-lab up -d} 后，NameNode、ResourceManager、NodeManager、Spark client 和 Spark History Server 能够启动；\texttt{start-dev.sh} 能在 \texttt{5051/5173} 启动后端与前端；HDFS 能保存样本数据；Spark submit 能产生 \texttt{SUCCEEDED} application；前端 Ops 页能显示 benchmark、HDFS 输入、YARN app id 和质量状态。

实验指标要求记录 elapsed seconds、cleaned rows/sec、input rows、output rows、Spark application status、driver peak memory、shuffle/spill 状态、failed/retried task count、output artifact row counts 和 quality gate status。其中最关键的结论阈值如下：

\begin{enumerate}
    \item YARN application 最终状态为 \texttt{SUCCEEDED}。
    \item driver 不再因为 full result \texttt{collect()} 或大 JSON payload OOM。
    \item recommendation 明细收集行数固定为 preview limit，质量统计由 Spark 聚合完成。
    \item \texttt{pair\_base\_rows / cleaned\_rows} 保持在可解释范围内，不能出现未受控的二次级组合膨胀。
    \item 5\% 样本相较 1\% 样本运行时间增长应接近线性，不能出现失控长尾。
    \item Spark History Server 能查看最终实验 event log；旧 event log 若触发 JSON 限制，应清理或用限制 plan string 后的日志替代。
\end{enumerate}

\section{数据清理需求}

旧数据集和旧实验日志会快速占用 Docker 与 HDFS 空间，因此系统需要及时清理冗余产物。清理策略是：保留原始 Kaggle Oct/Nov CSV；保留 1\% 与 5\% 本地样本；保留正式 benchmark 目录；保留最终 6 个 Spark History event log；删除早期 smoke、retry、history recheck benchmark 目录和旧的大 History 日志。清理后的存储快照见第 4 章图\ref{fig:storage-cleanup}。
